\documentclass[10pt,a4paper]{amsart}
\usepackage[margin=19mm]{geometry}
\usepackage{amsmath,amssymb,mathtools}
\usepackage[scr=rsfs,cal=euler]{mathalfa}
\usepackage{xfrac}
\usepackage[unicode,hidelinks,bookmarks=false]{hyperref}
\newcommand{\ie}{i.e.\ }
\newcommand{\cf}{cf.\ }
\newcommand{\resp}{respectively\ }
\newcommand{\Subsection}{\subsection}
\providecommand{\nobreakdash}{\nobreak\hbox{-}}
\setlength{\emergencystretch}{2em}
\multlinegap0pt
\usepackage{amssymb}
\newtheorem{definition}{Definition}
\newtheorem{lemma}{Lemma}
\newtheorem{proposition}{Proposition}
\newcommand{\A}{\mathcal{A}}
\newcommand{\B}{\mathcal{B}}
\newcommand{\CC}{\mathcal{C}}
\newcommand{\D}{\mathcal{D}}
\newcommand{\LL}{\mathcal{L}}
\newcommand{\N}{\mathbb{N}}
\newcommand{\Q}{\mathbb{Q}}
\newcommand{\ddd}{\boldsymbol{d}}
\newcommand{\dgsp}{\text{DgSp}}
\newcommand{\lom}{^{< \omega}} 
\title{Measuring the Complexity of Countable Presburger Models}
\author{Jason Block}
\date{Source: arXiv:2601.21118v3 (CC BY 4.0)}
\begin{document}
\maketitle

\noindent\emph{Source and licence.} Measuring the Complexity of Countable Presburger Models. Version arXiv:2601.21118v3 (CC BY 4.0), distributed by arXiv under the Creative Commons Attribution 4.0 International licence, http://creativecommons.org/licenses/by/4.0/, as recorded in the arXiv licence metadata for this exact version and shown on the abstract page https://arxiv.org/abs/2601.21118v3. Full licence notice: \texttt{licenses/arxiv-2601.21118-CC-BY-4.0.txt}.\par\smallskip

\noindent\textit{Editorial scope.} Counted unit: the article's proposition L (source0.tex lines 889--893). The article's complete original proof of it is delivered with it (source0.tex lines 895--957, one \texttt{proof} environment of 63 lines). No mathematics is written, repaired or re-derived: the statement and the proof are the article's own lines, in the article's own notation, and no retained line is substituted or re-flowed.  Retained uncounted as the source-local context the proof needs: lines 289--292; lines 867--870.  Nothing was omitted from any retained span, so the package carries no omission record.  No retained line contains a citation, so the wrapper carries no bibliography and no bibliography entry.  No retained line contains a figure environment, an image-inclusion command or drawing code, so no image is delivered and the package declares no asset dependency.  Editorial formatting only, in addition: an AMS \texttt{amsart} preamble that restates the article's own declarations and macros used by the retained lines, namely the environments \texttt{definition}, \texttt{lemma}, \texttt{proposition} on separate counters, together with the standard packages those lines need.  The retained environments print in this document as Definition 1, Lemma 1, Proposition 1 in order of appearance, whereas the article prints the same items with the numbering of its own full document.  The complete native source of the article as distributed in the arXiv e-print for this exact version is bundled unchanged as \texttt{sources/193572/source0.tex}.
\begin{definition}
\label{arch equiv}
Let $G$ be an ordered group (not necessarily divisible). Given $x\in G$, define $|x|=\max(x,-x)$. We say that $x,y\in G$ are Archimedean equivalent (denoted $x\sim y$) if there exists a positive integer $N$ such that $N|x|>|y|$ and $N|y|>|x|$. If $N|x|<|y|$ for all $N\in \N$, then we say $x\ll y$. 
\end{definition}

\begin{lemma}
\label{c.e. in}
Let $\A$ be a copy of a divisible ordered abelian group.  If the atomic diagram $\D(\A)$ is c.e.\ in some degree $\ddd$, then $\D(\A)$ is $\ddd$-computable.    
\end{lemma}

\begin{proposition}
\label{V_L  d' in L}
Given any countable linear order $\LL$,
    $\dgsp(V_\LL)=\{\ddd:\ddd'\in \dgsp(\LL)\}. $
\end{proposition}

\begin{proof}
    Let $\CC$ be a copy of $V_\LL$ of degree $\ddd$. The Archimedean equivalence relation $\sim$ (as in Definition \ref{arch equiv}) is c.e.\ in $\ddd$, and so $(\CC-\{0_\CC\})/\!\sim$, which is isomorphic to $\LL$ as a linear order, is $\ddd'$-computable. Hence, $\dgsp(V_\LL)\subseteq \{\ddd:\ddd'\in \dgsp(\LL)\}$.  

    For the other direction, suppose that $\ddd'\in \dgsp(\LL)$. It follows that there is some $\ddd$-computable set $T$ such that $T'=\{e:\Phi^T_e(e)\downarrow\}$ computes a copy $\A$ of $\LL$. Thus, there exists some Turing functional $\Psi$ such that $\Psi^{T'}(i,j)=1$ if $a_i<_\A a_j$ and $\Psi^{T'}(i,j)=0$ if $a_j<_\A a_i$. We use our oracle $T$ and functional $\Psi$ to compute the atomic diagram of a copy $\B$ of $V_\LL$. We do so by continually "guessing" at initial segments $\sigma$ of $T'$. We use the notation $\preceq$ and $\prec$ to denote initial segments and proper initial segments respectively.  Note, if we define $\sigma_{n,m}\in 2^n$ such that $\sigma_{n,m}(e)=1$ if and only if $\Phi^T_{e,m}(e)\downarrow$, we have that $\sigma_{n,m}$ will be an initial segment of $T'$ for sufficiently large $m$. 
    
    At each stage $s$ of the construction, we will have a set $H_s$ of elements in $\B$ of the form $h_{\sigma,i}$ for $i\in \N$ and $\sigma$ that still "look like" they may be initial segments of $T'$. That is, $\sigma$ such that $\sigma(e)=0$ if and only if $\Phi^T_{e,s}(e)\uparrow$. We treat all elements of $H_s$ as if they are in their own Archimedean equivalence class.  If at a later stage $t$ we have that $\Phi^T_{e,t}(e)\downarrow$ for an $e$ such that $\sigma(e)=0$, then we conclude that $\sigma$ is not an initial segment of $T'$ and we make it so each element of $H_t$ of the form $h_{\sigma, i}$ will actually be Archimedean equivalent to some other element of $H_t$ and we exclude $h_{\sigma,i}$ from $H_{t+1}$. Defining $H=\lim_sH_s$, we will have that all elements of $H$ are of the form $h_{\sigma,i}$ for $\sigma\prec T'$. Given $h_{\sigma,i}, h_{\rho,j}\in H$, we will have $h_{\sigma,i}\ll h_{\rho,j}$ if and only if $a_i<_\A a_j$. Additionally, each nonzero Archimedean equivalence class of $\B$ will contain exactly one element of $H$, and every element of $\B$ will be a finite $\Q$-linear combination of elements in $H$. Thus, $\B\cong \bigoplus_{i\in \A} \Q \cong V_\LL$. 

    For simplicity, we will view  $\B $ as purely relational. Furthermore, since $\B$ is divisible, the element $qx$ with $q\in \Q$ and $x\in \B_n$ is well defined and can be found computably in the atomic diagram $\D(\B)$.  
Thus, we may add multiplication by each $q\in \Q$ to our language. We get that the sentences in $\D(\B)$ will be those of the following forms:
\begin{itemize}
    \item $b_i=0$
    \item $b_i<b_j$
    \item $qb_i=b_j$
    \item $b_i + b_j = b_k$
\end{itemize}
where $b_i,b_j,b_k\in \B$ and $q\in \Q$.
We now give instructions for how, using our $\ddd$-oracle $T$, one can enumerate $\D(\B)$. Note, it will thus follow from Lemma \ref{c.e. in} that $\D(\B)$ is $\ddd$-computable.  \\

\paragraph{Construction:}  We let $I_s$ denote the set of possible initial segments of $T'$ through stage $s$. That is, $I_s=\{\sigma \in 2\lom: |\sigma|\leq s \text{ and } \Phi^T_{e,s}(e)\downarrow \text{ if and only if }\sigma(e)=1\}$. We make use of a function $f$ that will tell us for how large an $n\in \N$ the function $\Psi^\sigma$ can tell the order of all elements $a_i\in \A$ with $i<n$ when only running for $s$ steps. Specifically given $\sigma\in 2\lom$, we define $f(\sigma, s)$ to be the largest natural number $n\leq s$ such that $\Psi^\sigma(i,j)$ halts within $s$ steps for all $i,j<n$. Note, as $\sigma$ approaches $T'$ and $s$ approaches infinity, $f(\sigma,s)$ approaches infinity. 

For each $s>0$, define the set $X_s\subseteq I_t \times \N$ to be all elements of the form $(\sigma,m)$ that satisfy the following:

\begin{itemize}
    \item[(i)] $m\leq f(\sigma,s)$; and
    \item[(ii)] $m>f(\rho,s)$ for all $\rho\prec \sigma$.
\end{itemize}

\noindent Define $H_s=\{h_{\sigma,m}: (\sigma,m)\in X_s\}$. We will define an ordering on $H_s$ so that if $\rho\preceq \sigma$, then $h_{\rho,m}<h_{\sigma,n}$ if and only if $\Psi^\sigma (m,n)=1$. The purpose of condition (ii) is so $H=\lim H_s$ will contain exactly one element of the form $h_{\sigma, m}$ for each $m\in \N$ and so $H$ will be isomorphic to $\A$ as a linear order.

Let $\tau$ be the shortest initial segment of $T'$ such that $\Psi^\tau(0,1) \downarrow$. Let $r$ be the number of stages it takes for $\Psi^\tau(0,1)$ to halt. Define $t=\max(|\tau|,r)$. For reasons that will become clear later, we begin our construction at stage $t$ instead of stage $0$. \\

\subparagraph{Stage $t$:} Add the sentence $b_0=0$ to $\D(\B)$. For each $h \in H_t$, pick a fresh element $b_i \in \B$ and define $b_i=h$. Add the sentence $h>0$ to $\D(\B)$ for all $h\in H_t$. Given $h_{\sigma, m}, h_{\rho,n}\in H_t$, add the sentence $h_{\rho, m}< h_{\sigma,n}$ to $\D(\B)$ if either: 

\begin{itemize}
    \item $\rho \preceq \sigma $ and $\Psi^\sigma(m,n)=1$; or
    \item $\sigma \preceq \rho$ and $\Psi^\rho(m,n)=1$; or
    \item $\sigma$ and $\rho$ are incomparable but $\rho$ comes before $\sigma$ lexicographically.\\ 
\end{itemize}



\subparagraph{Stage $s+1>t$:}\,\\
\subparagraph{(i)} If $\sigma\in I_s - I_{s+1}$, then we will have that every $h_{\sigma,i}\in H_s$ will not be in $H_{s+1}$. Thus we arrange that all such $h_{\sigma,i}$ will be Archimedean equivalent to some $h\in H_{s+1}$. To do so, simply add the sentence $q h_{\sigma,i}= h$ for some sufficiently large (or small) $q\in \Q$ and appropriate $h\in H_{s+1}$ in such a way that respects the order on $H_s$. Note, our choice to begin the construction at stage $t$ guarantees that there is still at least one element in $H_{s+1}$. For every $r\in \Q$ such that the sentence $r h = a_c$ is in $\D(\B)$, add the sentences $rqh_{\sigma, i}=b_c$ and $\frac{1}{rq}b_c=h_{\sigma,i}$ to $\D(\B)$ as well. 

Given $h_{\sigma, m}, h_{\rho,n}\in H_{s+1}$ such that neither $h_{\sigma, m}< h_{\rho, n}$ nor $h_{\rho, n}<h_{\sigma, m}$ is in $\D(\B)$, add the sentence $h_{\rho, m}< h_{\sigma,n}$ to $\D(\B)$ if either: 

\begin{itemize}
    \item $\rho \preceq \sigma $ and $\Psi^\sigma(m,n)=1$; or
    \item $\sigma \preceq \rho$ and $\Psi^\rho(m,n)=1$; or
    \item $\sigma$ and $\rho$ are incomparable but $\rho$ comes before $\sigma$ lexicographically.\\ 
\end{itemize}

\subparagraph{(ii)} Let $(q_i)_{i\in \N}$ be an enumeration of the rationals. For each $h\in H_{s+1}$, take the least $i$ such that there is no sentence of the form $q_i h=b$ in $\D(\B)$. Take a fresh element $b_c$ and add the sentences $q_i h = b_c$ and $\frac{1}{q_i} b_c=h$ to $\D(\B)$. \\

\subparagraph{(iii)} For every finite tuple $(h_1,...,h_n)\in H\lom$ with each $h_i$ unique and every tuple $(q_1,...,q_n)\in\Q^n $ such that a sentence of the form $q_i h_i= b$ is in $\D(\B)$ for each $i$ but there is no sentence of the form $q_1h_1 + \cdots +q_nh_n= b $ in $\D(\B)$, take a fresh element $b_c$ and add the sentence $q_1h_1 + \cdots +q_nh_n= b_c$ to $\D(\B)$. \\

\subparagraph{(iv)} If we have $b,b'\in \B$ that have been previously mentioned in the construction but we have neither $b<b'$ nor $b'<b$ in $\D(\B)$, we must now decide the ordering on $b,b'$. Fix an enumeration $h_1,...,h_n$ of all the elements of $H$. We have that there exists $q_1,...,q_n,q'_1,...q'_n\in \Q$ such that $b=q_nh_n+ \cdots +q_1 h_1$ and $b'= q'_nh_n + \cdots +q'_1 h_1$. Add the sentence $b<b'$ to $\D(\B)$ if and only if $(q_n,...,q_1)$ is lexicographically less than $(q'_n,...,q'_1)$. 

Last, add $b_i+ 0 = b_i$, $1b_i=b_i$ and $0b_i= 0$ to $\D(\B)$ if they are not already in. If $b_i+b_j=b_c$ is in $\D(\B)$, add $b_j + b_i = b_c$ as well if it is not already in. If $qb_i=b_j$ is in $\D(\B)$ with $q\neq 0$ add $\frac{1}{q}b_j=b_i$ as well if it is not already in. This completes the construction.\\

\paragraph{Verification}
Define $H=\lim_sH_s$. Note that $h_{\sigma,n}\in H$ only if $\sigma\prec T'$. Furthermore, if $h_{\rho,m}\in H$ with $\rho\preceq\sigma$, then $h_{\rho,m}<_\B h_{\sigma,n}$ if and only if $\Psi^\tau(m,n)=1$ if and only if $\Psi^{T'}(m,n)=1$ if and only if $a_m<_\A a_n$. Hence, $(H,<_\B)$ is isomorphic to $\LL$. Additionally, our instructions give that $H$ is exactly the Archimedean rank of $\B$, and every element of $\B$ has the form $q_1 h_1 + \cdots + q_nh_n$ with each $q_i\in \Q$ and $h_i\in H$. Hence, $\B\cong \bigoplus_{h\in H} \Q \cong \bigoplus_{l\in \LL} \Q = V_\LL$. 
\end{proof}

\end{document}
